\documentclass[12pt]{article}
\usepackage[utf8]{inputenc}
\usepackage[T1]{fontenc}
\usepackage{hyperref}
\title{Mixtec, Cuyamecalco Language Resources for K-12 Curriculum Development}
\author{Compiled by Claude Sonnet 5.5 (claude-sonnet-5-5) for Andr\'as Kornai}
\date{}
\begin{document}
\maketitle

\section{Introduction}
Cuyamecalco Mixtec (ISO 639-3 xtu; Glottocode cuya1240) is an Otomanguean language of Oaxaca, Mexico. It is spoken in Cuyamecalco, San Miguel Santa Flor and Santa Ana Cuauhtémoc. Estimates are about 1,900 speakers (a 2000 count), with roughly 70 monolinguals in San Miguel. Wikitongues gives the endonym Tu'un savi and rates it EGIDS 6b (threatened). The search found only catalogue and reference entries, and no dictionary, grammar, primer or textbook.

\section{Primary Educational Resources}
No primers, textbooks or school programmes were found.

\section{Linguistic and Reference Materials}
\begin{itemize}
\item Wikipedia, Cuyamecalco Mixtec: \url{https://en.wikipedia.org/wiki/Cuyamecalco_Mixtec}
\item Endangered Languages Project entry (listed under the name Northern Alta Mixtec): \url{https://www.endangeredlanguages.com/lang/10533?hl=fr}
\item Wikitongues language page: \url{https://wikitongues.org/languages/xtu/}
\end{itemize}

\section{Religious and Community Resources}
Joshua Project reports no written standard, no Bible translation and no Jesus Film or audio Bible for this language: \url{https://joshuaproject.net/languages/xtu}

\section{Digital and Audio Resources}
Global Recordings Network has a page for this language, but the search found no audio recordings listed on it: \url{https://www.globalrecordings.net/en/language/xtu}

\section{Curriculum Development Context}
The language is threatened and has no known written standard or teaching materials. A curriculum would need new documentation and orthography work with the community. Nothing found here is ready for classroom use.
\end{document}